\chapter{Conclusão}\label{cap:conclusao}
\addcontentsline{toc}{chapter}{Conclusão}

Este capítulo encerra o trabalho retomando, à luz dos resultados discutidos no Capítulo~\ref{cap:avaliacao}, o problema de pesquisa e as questões formuladas na Seção~\ref{sec:questoes-pesquisa-metodologia}. A Seção~\ref{sec:sintese-trabalho} sintetiza o percurso do trabalho. A Seção~\ref{sec:resposta-questoes-pesquisa} responde diretamente a cada questão de pesquisa. A Seção~\ref{sec:contribuicoes} indica as contribuições do estudo. A Seção~\ref{sec:limitacoes-conclusao} consolida as limitações reconhecidas. A Seção~\ref{sec:trabalhos-futuros} aponta direções de continuidade.

\begin{quote}
\textbf{Nota de redação:} como nos capítulos anteriores, os trechos entre colchetes marcam pontos que só podem ser redigidos após a coleta e a análise dos resultados descritos no Capítulo~\ref{cap:avaliacao}. Nenhuma conclusão factual foi antecipada nesta versão.
\end{quote}

\section{Síntese do Trabalho}\label{sec:sintese-trabalho}

[A PREENCHER APÓS A COLETA: retomada breve do problema de pesquisa --- comparação entre SAST tradicional, agentes empresariais de IA e abordagem híbrida SAST + IA na detecção de vulnerabilidades em Python --- e do percurso metodológico e experimental adotado, sem repetir integralmente o conteúdo dos capítulos anteriores.]

\section{Resposta às Questões de Pesquisa}\label{sec:resposta-questoes-pesquisa}

[A PREENCHER APÓS A COLETA: resposta objetiva a cada uma das questões QP1--QP5 definidas na Seção~\ref{sec:questoes-pesquisa-metodologia}, com remissão direta às seções correspondentes do Capítulo~\ref{cap:avaliacao} onde cada resultado foi apresentado.]

\begin{itemize}
\item \textbf{QP1} --- desempenho comparativo de detecção: [A PREENCHER, ver Seção~\ref{sec:desempenho-quantitativo}].
\item \textbf{QP2} --- qualidade dos relatórios verdadeiros-positivos: [A PREENCHER, ver Seção~\ref{sec:qualidade-relatorios}].
\item \textbf{QP3} --- efeito de fornecer alertas SAST aos agentes de IA: [A PREENCHER, ver Seção~\ref{sec:efeito-hibrido}].
\item \textbf{QP4} --- estabilidade entre execuções e adequação da priorização: [A PREENCHER, ver Seções~\ref{sec:estabilidade-execucoes} e~\ref{sec:priorizacao-resultados}].
\item \textbf{QP5} --- tempo de execução e custo observável: [A PREENCHER, ver Seção~\ref{sec:tempo-custo-resultados}].
\end{itemize}

\section{Contribuições}\label{sec:contribuicoes}

[A PREENCHER APÓS A COLETA: contribuições efetivamente entregues pelo trabalho, por exemplo a comparação tripla SAST--IA--híbrido em Python identificada como lacuna na Seção~\ref{sec:lacunas-literatura}, o \textit{harness} e os artefatos de reprodução descritos no Capítulo~\ref{cap:desenvolvimento}, e a rubrica qualitativa aplicada aos relatórios das seis condições.]

\section{Limitações}\label{sec:limitacoes-conclusao}

[A PREENCHER APÓS A COLETA: consolidação das limitações já antecipadas na Seção~\ref{sec:validade-limitacoes} e das limitações adicionais identificadas apenas durante a execução, registradas na Seção~\ref{sec:ameacas-validade-observadas}, incluindo a delimitação da generalização das conclusões ao RealVuln v1.0 e às versões, regras, \textit{prompts} e ferramentas efetivamente avaliadas.]

\section{Trabalhos Futuros}\label{sec:trabalhos-futuros}

[A PREENCHER APÓS A COLETA: direções de continuidade sugeridas pelos resultados obtidos --- por exemplo, extensão a outros benchmarks ou linguagens, avaliação de outras ferramentas empresariais de IA, refinamento dos desenhos híbridos discutidos na Seção~\ref{sec:abordagens-hibridas}, ou aprofundamento da avaliação de priorização e severidade.]
